\section{Related Work}\label{sec:related}

\subsection{Detector architectures and benchmark design}

RawNet2 learns representations straight from speech waveforms whereas AASIST learns with graph attention applied to spectro-temporal representations \cite{rawnet,aasist}. A compact convolutional backbone represents another trade-off between compute and representations. MobileNetV3 was developed for visual recognition on mobile devices \cite{mobile}. In our case, this architecture encodes consecutive Mel-based feature sequences, which are aggregated in time using a recurrent head. There exist published lightweight approaches that combine a distinct feature extraction scheme and contrastive learning \cite{din}. The previous detector designs inform this approach, and the controlled intervention is carried out with a consistent backbone and head in all arms.

ASVspoof 2019 separates logical access, synthetic, or converted speech from physical access, replayed samples \cite{asv19}. The ASVspoof 2021 campaign expands benchmarking to cover distinct logical access, physical access, and deepfake challenges under distinct conditions \cite{asv21}. ASVspoof 5 scales the crowdsourced data collection and spoofs and adversary testing \cite{asv5}. In-the-Wild data illustrate why a split in the benchmark dataset alone does not guarantee transport to collected in-the-wild data \cite{wild}. The current development dataset contains separated speakers from training as well as attack IDs present in training. Within the release, in terms of algorithms, A16 and A19 have commonalities with A04 and A06, respectively, and A17 with A05 \cite{asv19}. Such relations place certain limitations on how held-out attack IDs can be interpreted.

The multilingual MLAAD dataset targets variations in speech synthesis systems and languages \cite{mlaad}. This dataset design offers an orthogonal test of transport to languages and generation domains.

\subsection{Signal transformations and robust learning}

Opus is a standard lossy codec \cite{opus}. Compression may affect spectral and temporal features through content-dependent coding choices. Resynthesis using a learned neural codec is another kind of transformation family, with artifacts of its own. The discussion and case study by Maltby et al. list codecs along with dataset bias and other failure modes \cite{maltby}. Such work motivates a separation of operator, representation, and detector.

RawBoost mixes convolutive, impulsive, and stationary waveform distortions and demonstrates improvements on ASVspoof 2021 LA \cite{rawboost}. The current augmentation selects only one out of the gains, noise, codec, and clean views for each example and applies binary cross-entropy loss. Frequency-selective adversarial training modifies the selection criterion for the perturbations \cite{fsat}. Group distributionally robust optimization assigns weights based on the risk of the group \cite{dro}. DK-CAST provides additional supervision by learning clean-to-distorted distillation \cite{dkcast}. SONAR examines spectral-contrastive residual representations across a wide range of shifts \cite{sonar}. These methods tackle objective or representation changes. The current experiment keeps the representation family and the objective fixed to measure the effect of the view policy.

\begin{table*}[t]
\caption{Comparison of research designs. Different cohorts and operating rules prevent a direct performance ranking.}\label{tab:prior}
\centering
\begin{tabular}{@{}p{.15\textwidth}p{.22\textwidth}p{.25\textwidth}p{.27\textwidth}@{}}
\toprule
Study & Data or evaluation setting & Treatment or model & Relationship to the present comparison\\
\colrule
RawBoost and F-SAT \cite{rawboost,fsat} & ASVspoof and DeepFakeVox-HQ, respectively & Waveform augmentation and frequency-selective adversarial training & Establish robust training precedents. Neither is retrained as a matched arm here.\\
DK-CAST and SONAR \cite{dkcast,sonar} & Speech spoofing and domain-shift evaluations & Distillation and spectral residuals & Intervene on learning objectives or representations. The present study intervenes on training views.\\
ADD-C and bona fide cross-testing \cite{addc,kwok} & Simulated communication and diverse authentic speech & Codec and loss grids, and bona fide domain testing & Motivate stress coverage and authentic-speech reporting. The present study controls order on a fixed source.\\
Present study & Three development initialization blocks, shared development and locked within-release held-out & Matched lightweight detector, ordinary augmentation & Links ordered processing and class-specific error transitions on the same source. Broader transport and mechanisms remain open.\\
\botrule
\end{tabular}
\end{table*}

\subsection{Order effects and operating-point uncertainty}

ADD-C tests detection under conventional codecs and simulated packet loss and evaluates augmentation for those communication conditions \cite{addc}. Bona fide cross-testing varies authentic-speech domains against spoofed sets and shows why authentic-speech coverage matters \cite{kwok}. Both motivate class- and condition-specific evaluation. Here the comparison reverses noise and conventional Opus on the same source with a shared noise plan, and measures how a matched training-view intervention changes the resulting threshold errors. This narrower intervention yields source-level evidence about the two orders under the tested settings.

Calibration concerns the interpretation of scores, while selective prediction introduces a reject action \cite{guo,selective}. Work on calibrated audio deepfake detection illustrates the importance of these distinctions in the speech domain \cite{pascu}. In this case, the sigmoid prediction is represented by p, which serves as an uncalibrated score. Rank improvement along with error increment at the same threshold is feasible. The following table (Table~\ref{tab:prior}) illustrates the way the closest prior research differs from this one in terms of scope.

\subsection{Channel exposure and representation choices}

Work with a pretrained wav2vec 2.0 front end and data augmentation provides published evidence for a different route to addressing channel and source shift \cite{wav2vecaugment}. The present comparison uses a compact Mel-sequence classifier with the same representation family in both arms. This limits the treatment comparison to exposure during fitting and avoids interpreting a change of backbone, input bandwidth, or pretraining corpus as an augmentation effect. The design also avoids comparisons with a published EER obtained under a different source partition and operating protocol.

Cross-domain multilingual benchmarks such as XMAD-Bench address variation across synthesis and language domains \cite{xmad}. Such benchmarks answer a different question from a source-matched transformation test. Pairing ensures that the source is held constant throughout the processing of the channel, which aids in making the contrast of the channel more interpretable. Pairing also makes sure that the outcome is conditioned on the population of the source within the chosen cohort. The cross-domain benchmark allows a larger population but creates extra differences that need to be taken into account in a score change.

Group-robust learning prioritizes a particular group criterion and usually demands the selection of groups and weighting. On the contrary, the conventional augmentation optimization minimizes the average loss over samples. It is important due to the fact that the latter objective does not have to secure the worst subgroup even when it enhances discrimination as a whole. The journal provides the comparison of the traditional policy intervention and subgroup effect measurement. No class-dependent weighting, consistency loss, and group-robust optimization are proposed.

The previous schemes also differ from each other in terms of the statistical problem posed by the operating rule. EER computes an equal error estimate of the score distribution, while a fixed threshold allows estimating separately both false alarms and misses for each processing configuration. Therefore, we retain the registered threshold and provide the analysis of the curve near it. The attack and speaker differences depend on the model fit. They supplement the benchmark comparison without making assumptions about a different source population, intervention, or operating rule.
